\section{The matrix-game gap}\label{sec:matrix-gap}

The matrix game's soundness follows by comparing two bounds for the same
advice experiment.  Proposition~\ref{prop:matrix-decoding} extracts local
strategies with a fixed positive agreement probability from test acceptance
at least $0.99$.  Lemma~\ref{lem:zero-advice} makes every such agreement probability
tend to zero on a NO parity instance.  We state the inputs and give the
deduction here; their proofs occupy Sections~\ref{sec:matrix}
and~\ref{sec:zero-advice}.

\subsection{The shortcode soundness theorem}

The structural input describes functions that often keep their value
under a uniform rank-one perturbation.  Such a function agrees with an
affine evaluation on a set obtained by fixing a bounded number of linear
combinations of rows and columns.

Write $\operatorname{Mat}_{\ell,m}=\Hom(\F_2^m,\F_2^\ell)$.
A \emph{row and column slice} is a nonempty set specified by equations
\begin{equation}
 \mathcal S=
 \{M:d_i^\top M=s_i^\top\ (1\le i\le n_r),\quad
          Mq_j=t_j\ (1\le j\le n_c)\},
 \label{eq:shortcode-slice}
\end{equation}
where $d_i\in\F_2^\ell$, $s_i\in\F_2^m$, $q_j\in\F_2^m$, and
$t_j\in\F_2^\ell$.  The counts $n_r,n_c$ need not be the ranks of their
respective systems.

\begin{theorem}[Equality soundness for degree-two shortcode]
\label{thm:shortcode}
For every $\eta\in(0,1)$ there are $\alpha\in(0,1]$, an integer
$r_{\rm s}\ge1$, and an integer $\ell_0$, depending only on $\eta$,
with the following property.  For every $\ell\ge\ell_0$ there is
$m_0=m_0(\eta,\ell)$ such that, for every $m\ge m_0$ and every function
$f:\operatorname{Mat}_{\ell,m}\to\F_2^\ell$, if
\begin{equation}
 \Pr_{M,a,l}\bigl[f(M)=f(M+a l^\top)\bigr]\ge\eta,
 \label{eq:shortcode-acceptance}
\end{equation}
then there are a nonempty row and column slice $\mathcal S$ as in
Equation~\eqref{eq:shortcode-slice} with
$n_r,n_c\le r_{\rm s}$, $z\in\F_2^m$, and $u\in\F_2^\ell$ such that
\begin{equation}
 \Pr_{M\mid\mathcal S}[f(M)=Mz+u]\ge\alpha.
 \label{eq:shortcode-conclusion}
\end{equation}
Here $M\in\operatorname{Mat}_{\ell,m}$, $a\in\F_2^\ell$, and $l\in\F_2^m$ are independent and
uniform, and $M\mid\mathcal S$ is uniform on the slice.  No folding
condition is imposed on $f$.
\end{theorem}


Theorem~\ref{thm:shortcode} is the soundness statement of Barak,
Kothari, and Steurer for the equality form of their degree-two shortcode
test~\cite[Hypothesis~2.1 and Section~1.6]{BKS18}, with the dimension
requirements stated explicitly: first $\ell\ge\ell_0$, then
$m\ge m_0(\eta,\ell)$.  They derive it from their hypothesis on
non-expanding sets of the shortcode graph by passing to a fiber of
$f$~\cite[Hypothesis~1.8 and Lemma~2.3]{BKS18}, and show that this
hypothesis is equivalent to the Grassmann expansion
hypothesis~\cite[Theorem~1.11]{BKS18}, which Khot, Minzer, and Safra
proved~\cite[Theorem~1.12]{KMS23}.  The proof in
Section~\ref{subsec:shortcode-proof} gives this derivation directly: it
applies Grassmann expansion to a fiber of $f$ in the matrix chart.

The constants and dimension thresholds in Theorem~\ref{thm:shortcode}
can be chosen effectively, with rational $\alpha$ when $\eta$ is rational.
This follows from the quantitative bound in
\cite[Theorem~2.13 and Section~2.4]{KMS23} and the comparison estimates
in \cite[Lemmas~2.7--2.8]{KMS23}, with theorem numbers referring to
the 2018 revision.  The proof in Section~\ref{subsec:shortcode-proof}
gives explicit sufficient inequalities for these choices.

We now fix the parameters needed below.  Choose $p_*>0$ and the associated
$r_*$ from Lemma~\ref{lem:equivariant-noise}, and put $\eta=2^{-r_*}/16$.
For this $\eta$, take $\alpha,r_{\rm s},\ell_0$ from
Theorem~\ref{thm:shortcode}.  Then choose $\ell\ge\ell_0$ sufficiently
large for Lemma~\ref{lem:equivariant-noise} and for
\begin{equation}
 2^{-(\ell-r_{\rm s})}<\alpha/8,\qquad 2^{-\ell}<\eta/4.
 \label{eq:matrix-parameters}
\end{equation}
The value of $\eta$ is
the acceptance level at which Theorem~\ref{thm:shortcode} is applied:
Lemma~\ref{lem:matrix-spectral} gives average ordinary acceptance at least
$0.9\cdot2^{-r_*}=14.4\eta$, which leaves room for the erasure in
Lemma~\ref{lem:row-advice-agreement}.  The first inequality in
Equation~\eqref{eq:matrix-parameters} supplies slack for erasure and
folding in Lemma~\ref{lem:row-advice-agreement} and for the Fourier estimate in
Equation~\eqref{eq:nontrivial-correlation}.

Obtain the map and noise data $(\mathcal V,K,C,\mu)$ at this value of $\ell$,
and fix a linear splitting $\mathcal V=K\oplus W$.  All these objects
are fixed before $k$ grows.  In particular, $\ell>r_{\rm s}$.
We write a table as $P=(M,T)$, with $M:L_U\to K$ and $T:L_U\to W$.

\subsection{The advice experiment and the decoding statement}

We first specify the local inputs for two strategies whose answer agreement
will contradict the outer-game gap.  This experiment is defined for every positive integer $k$
and every $0<\beta<1$.  Let $\mathcal Q_k$ consist of all equation
questions $U\in\mathcal E^k$ and all questions $O$
obtained from them by the coordinate projections of Section~\ref{sec:game}.
The questions $U,O$ and the projection
$\bar\pi:L_U\to L_O$ have the law in Section~\ref{sec:game}: independently
in each block, an equation is retained with probability $1-\beta$, and a
uniformly chosen variable position is retained with probability $\beta$.

\begin{definition}[Projected advice experiment]
\label{def:projected-advice}
Sample $U,O,\pi$, then an independent uniform map
$A:K\to\F_2^{r_{\rm s}}$ and independent uniform maps
\[
 T_O:L_O\to W,\qquad M_O:L_O\to K.
\]
Put $S_O=AM_O$ and send the following inputs to the two provers:
\begin{equation}
 \begin{aligned}
  &\text{$U$-prover:}&& (U,A,T_O\bar\pi,S_O\bar\pi),\\
  &\text{$O$-prover:}&& (O,A,T_O,S_O).
 \end{aligned}
 \label{eq:projected-advice}
\end{equation}
The map $M_O$ beyond the stated advice is hidden.  Each prover returns an
actual answer, that is, a vector of first bit one in its coordinate
space.  They succeed when their answers $a_U,a_O$ satisfy
$\bar\pi a_U=a_O$.
\end{definition}

The $U$-message contains the maps
$(T_U,S_U)=(T_O\bar\pi,S_O\bar\pi)$, obtained by padding maps sampled
on $O$.  We will first find advice admitting agreement on a slice under a different law: an
unrestricted uniform table $(M,T)$ on $U$, with $S_U=AM$.  Later we
compare the full padded table $(M_O\bar\pi,T_O\bar\pi)$ with this
unrestricted table after averaging over the projection choices.  This
comparison is the bridge from the extraction argument to the experiment.

\begin{proposition}[Matrix decoding]
\label{prop:matrix-decoding}
Fix $p_*>0$ and choose $r_*$, $\eta$, $\alpha$, $r_{\rm s}$, $\ell_0$, $\ell$,
the map and noise data and the splitting $\mathcal V=K\oplus W$, in this order,
as in the paragraph containing Equation~\eqref{eq:matrix-parameters};
all of them are fixed before $k$.  There are a constant $\gamma>0$ and an integer $k_0\ge8$, depending only on
these fixed parameters, with the following property.

For every integer cube $k\ge k_0$ and every weighted parity instance
with three distinct variable IDs per equation, suppose a function
\[
 F_Q:\Hom(L_Q,\mathcal V)\longrightarrow K
\]
is given for every $Q\in\mathcal Q_k$.  Assume folding for every such
$Q$, every $P\in\Hom(L_Q,\mathcal V)$, and every $c\in K$, and assume
sharing for every coordinate projection $\bar\pi:L_U\to L_O$ and every
$P'\in\Hom(L_O,\mathcal V)$:
\begin{equation}
 \begin{aligned}
 F_Q(P+c e_0^\top)&=F_Q(P)+c,\\
 F_U(P'\bar\pi)&=F_O(P').
 \end{aligned}
 \label{eq:matrix-family-identities}
\end{equation}
If this family passes the table test of Equation~\eqref{eq:matrix-test}
with probability at least $0.99$, then the projected advice experiment
with $\beta=k^{-2/3}$ has local strategies whose probability of full
answer agreement is at least $\gamma$.
\end{proposition}

Here, and in Lemma~\ref{lem:zero-advice} and Proposition~\ref{prop:matrix-gap},
$k$ ranges over integer cubes so that $\beta=k^{-2/3}$ is rational; the
search for $k$ in Section~\ref{subsec:effective-k} uses this to check its
conditions exactly.

Every labeling of the vertices of $\mathcal G$ supplies such a family:
its restored functions satisfy Equation~\eqref{eq:matrix-family-identities}
by Equations~\eqref{eq:table-folding} and~\eqref{eq:table-sharing}.  The family is fixed when the
two response rules are defined and may be built into those rules.


The proof of Proposition~\ref{prop:matrix-decoding} is in
Section~\ref{sec:matrix}.  Its constant is uniform in the instance and
in the table family, which allows the following upper bound to contradict it.

\subsection{The upper bound and the gap}

The upper bound uses the same fixed advice spaces as
Definition~\ref{def:projected-advice}; it applies to every pair of local
strategies, independently of how they were obtained.

\begin{lemma}[Agreement bound from vanishing advice slopes]
\label{lem:zero-advice}
Let $\mathcal E$ be a weighted parity instance with three distinct
variable IDs in each equation and $\OPT(\mathcal E)\le4/5$.
For every integer cube $k>1$, every pair of local randomized strategies
in Definition~\ref{def:projected-advice}, with $\beta=k^{-2/3}$, satisfies
\begin{equation}
 \Pr\bigl[\bar\pi(a_U)=a_O\bigr]
 \le \left(1-\frac{\beta\,2^{-\dim W-r_{\rm s}}}{3600}\right)^k
 \le
 \exp\!\left(
 -\frac{2^{-\dim W-r_{\rm s}}}{3600}\,k^{1/3}
 \right).
 \label{eq:zero-advice-agreement-bound}
\end{equation}
In particular, the upper bound tends to zero uniformly over the parity
instance as $k$ tends to infinity along cubes, with the advice spaces
fixed.
\end{lemma}

Section~\ref{sec:zero-advice} proves this bound by restricting to singleton
coordinates on which the advice slopes vanish, and applying
Theorem~\ref{thm:repetition} to the remaining equation--variable questions.

\subsection{Deduction of the matrix-game gap}

We can now compare the upper bound for every advice strategy with the
strategies extracted from a highly accepting matrix labeling.

\begin{proposition}[A fixed gap with adjustable completeness]
\label{prop:matrix-gap}
Fix $p_*>0$ and choose $r_*$, $\eta$, $\alpha$, $r_{\rm s}$, $\ell_0$, $\ell$,
the map and noise data and the splitting $\mathcal V=K\oplus W$, in this order,
as in the paragraph containing Equation~\eqref{eq:matrix-parameters}.  There is a threshold, depending only on these fixed data, such that for
every integer cube $k$ above it and every weighted parity instance
$\mathcal E$ with three distinct variable IDs per equation, the matrix
game $\mathcal G$ of Equation~\eqref{eq:matrix-test} satisfies
\[
 \begin{aligned}
 \OPT(\mathcal E)\ge1-\xi
 &\quad\Longrightarrow\quad
 \val(\mathcal G)\ge1-k\xi-p_*/2,\\
 \OPT(\mathcal E)\le4/5
 &\quad\Longrightarrow\quad
 \val(\mathcal G)\le99/100.
 \end{aligned}
\]
The threshold is also independent of $\xi$.
\end{proposition}

\begin{proof}
Completeness is Lemma~\ref{lem:outer-completeness}.  For soundness,
fix a labeling of $\mathcal G$.  Its restored table functions satisfy
the folding and sharing identities of Section~\ref{sec:game}.
Proposition~\ref{prop:matrix-decoding} supplies a constant $\gamma>0$
and, for every sufficiently large cube $k$, gives local advice
strategies with agreement at least $\gamma$ whenever those functions
have matrix-test acceptance at least $0.99$.  For all sufficiently large
cubes $k$, the upper
bound in Equation~\eqref{eq:zero-advice-agreement-bound} is smaller than
$\gamma$.
Lemma~\ref{lem:zero-advice} contradicts the existence of those strategies
whenever $\OPT(\mathcal E)\le4/5$.  All thresholds are uniform over
the parity instance.
\end{proof}

The order matters: $p_*$ can be made arbitrarily small before the
map and noise parameters are fixed, and $\xi$ can be made arbitrarily small
after $k$ is fixed.  Thus the same NO gap is compatible with any
prescribed small completeness error.
